\section{E. The discriminant over an arbitrary base}
\label{sec:geometry-E}

We use the quotient convention $\mathbf P(E)=\operatorname{Proj}(\operatorname{Sym}E)$.
Put $N=r+1$, and let $d\geq2$ be the degree (the letter $k$ in cote 23/006).
The manuscript parametrizes forms by a projective bundle. Thus its equation is
fibrewise nonzero. This condition matters over a general base.

\begin{theorem}[Known: integral normalization]\label{geo-E-normalization}
Let $F=\sum_{|\alpha|=d}c_\alpha X^\alpha$ over
$U=\mathbf Z[c_\alpha]$, in $N\geq1$ variables. With the resultant normalized by
$\operatorname{Res}(X_1^{e_1},\ldots,X_N^{e_N})=1$, put
\[
 a(N,d)=\frac{(d-1)^N-(-1)^N}{d},\qquad
 D=N(d-1)^{N-1},\qquad w=d(d-1)^{N-1}.
\]
There is a primitive integral polynomial $\operatorname{Disc}_{N,d}\in U$ such that
\[
 \operatorname{Res}(\partial_1F,\ldots,\partial_NF)
       =d^{a(N,d)}\operatorname{Disc}_{N,d}(F).
\]
It has degree $D$ in the coefficients, and
$\operatorname{Disc}_{N,d}(F\circ g)=\det(g)^w\operatorname{Disc}_{N,d}(F)$.
Over any field, its vanishing is equivalent to the existence, after a finite field
extension, of a nonzero common zero of $F,\partial_1F,\ldots,\partial_NF$.
\end{theorem}

The exact source is Demazure~\cite{Demazure2012}, \S5, Lemma 11 (content),
Definition 4 (normalization), Proposition 11(c),(e) (degree and covariance), and
Proposition 12 (critical zeros). These statements were opened. Division takes
place in $\mathbf Z[c_\alpha]$ \emph{before} specialization; it is never division
by an element of the eventual base ring. The primitive polynomial is fixed by
this positive-content convention. An alternative convention changes its sign.

\begin{theorem}[Proved from the integral theorem: the relative statement]
\label{geo-E-relative}
Let $S$ be any scheme, $E$ locally free of rank $N$, and
$M\hookrightarrow\operatorname{Sym}^d E$ an invertible subbundle. Let
$X\subset\mathbf P(E)$ be the corresponding relative degree-$d$ hypersurface.
There is a canonical section
\[
 \delta_X\in
 \Gamma\bigl(S,M^{-D}\otimes(\det E)^{\otimes w}\bigr).
\]
Its formation commutes with arbitrary base change, and $X\to S$ is smooth
if and only if $\delta_X$ generates this line bundle at every point of $S$.
In particular the exponent $(d-1)(r+1)$ in cote 23/006 is false in general;
the corrected exponent is $(r+1)(d-1)^r$.
\end{theorem}
\begin{proof}
On an open set trivializing $E$ and $M$, write the image of the chosen generator
of $M$ as a homogeneous polynomial $f$. Evaluate the integral polynomial of
Theorem~\ref{geo-E-normalization} on its coefficients. Multiplying the generator
of $M$ by $u$ multiplies this value by $u^D$. Changing the frame of $E$ by $g$
multiplies the value by $\det(g)^{-w}$, because coordinates in the new frame
are obtained by the inverse change of variables. These are exactly the
coefficient changes for a section of $M^{-D}\otimes(\det E)^w$. They give
compatible sections on overlaps; polynomial evaluation also proves arbitrary
base-change compatibility.

For a point $s\in S$, the subbundle hypothesis makes $f_s$ nonzero. Over an
algebraic closure of $\kappa(s)$, a projective hypersurface defined by this
nonzero form is singular precisely when $f_s$ and all its homogeneous partial
derivatives vanish at a projective point. This follows on each affine chart
from the hypersurface Jacobian criterion; Euler's identity ensures the
equivalence of the homogeneous and chart versions, without dividing by $d$
(on $X_i=1$, the missing homogeneous derivative is a linear combination of
the others and $df$). The integral theorem therefore says that $\delta_X(s)$
is nonzero precisely when this geometric fibre is smooth.

No separate fibrewise flatness criterion is needed: at a point with one
invertible chart derivative the standard Jacobian presentation makes $X/S$
smooth directly. Conversely a smooth
morphism has smooth geometric fibres. Thus the desired equivalence follows
point by point. Finally, a section of a line bundle generates it exactly when
it is nonzero in every residue field.
\end{proof}

\begin{example}[Proved and computed: normalization at the prime $2$]
\label{geo-E-conic}
For $f=ax^2+by^2+cz^2+dxy+exz+f_0yz$ the divided discriminant is
\[
 4abc-af_0^2-be^2-cd^2+def_0.
\]
It is one half of the integral Hessian determinant, not the Hessian
determinant itself. Over $\mathbf F_2$, $x^2+yz$ defines a smooth conic:
its derivatives are $0,z,y$, and a critical projective point would have
$y=z=0$, where the equation forces $x=0$. Its divided discriminant is $1$,
while the naive resultant of derivatives is $0$. The program
\texttt{code/geometry\_discriminant.py} computes the integral formula and
checks that the singular-locus ideal on each of the three standard affine
charts has Gr\"obner basis $\{1\}$ over $\mathbf F_2$. Thus it checks emptiness
over every field extension, not merely absence of $\mathbf F_2$-points.
\end{example}

\begin{example}[Refuted degree and proved elliptic consistency]
\label{geo-E-cubic}
For $\lambda(x^3+y^3+z^3)$ over $\mathbf Q[\lambda]$, the resultant of the
three partial derivatives is $3^{12}\lambda^{12}$, and the content divisor
is $3^3$. The divided discriminant is $3^9\lambda^{12}$, contradicting
the proposed coefficient degree $6$ already on this one-parameter family.
In cote 23/005, $\det E\simeq L^5$ and $M\simeq L^6$, so the corrected
discriminant line is
\[
 M^{-12}\otimes(\det E)^{12}\simeq L^{-12}.
\]
This Hodge-line identification can be checked independently. The relative
Euler sequence on the quotient projective bundle gives
$\omega_{\mathbf P(E)/S}=\mathcal O(-3)\otimes\pi^*\det E$.
The cubic conormal is $\mathcal O_X(-3)\otimes p^*M$, so taking
determinants of its conormal sequence gives
$\omega_{X/S}=p^*(\det E\otimes M^{-1})=p^*L^{-1}$ on the smooth
locus in the base. Pushing forward
$0\to\pi^*M\otimes\mathcal O(-3)\to\mathcal O\to\mathcal O_X\to0$
gives $p_*\mathcal O_X=\mathcal O_S$: the zeroth and first direct images
of $\mathcal O(-3)$ on $\mathbf P^2_S$ vanish, as the standard three-chart
\v Cech calculation shows over any base ring. Hence the Hodge line
$p_*\omega_{X/S}$ is $L^{-1}$ by the projection formula, and the
discriminant line above is its twelfth tensor power. Both the coefficient
degree and the determinant weight are needed for this check.
\end{example}

If one drops the fibrewise-nonzero condition, the assertion as a statement
about arbitrary equations is false: the zero equation defines all of
$\mathbf P^{N-1}_S$, which is smooth, but has zero discriminant. For degree
$1$, a fibrewise-nonzero linear form defines a smooth relative hyperplane;
there is no discriminant divisor and the relevant invariant is $1$.
